\clearpage
\VerdierPhysicalPageAnchor{0266}%
\begingroup
\small
\setlength{\parindent}{0pt}%
\setlength{\parskip}{0pt}%
\setlength{\emergencystretch}{2em}%
\newcommand{\VerdierBibRow}[3]{%
  \VerdierSemanticAnchor{#1}%
  \par\noindent\hangindent=3em\hangafter=1
  \makebox[3em][l]{[#2]}#3\par\smallskip}

\vspace*{4.5\baselineskip}
\VerdierSemanticAnchor{bib-title}%
\begin{center}
  {\Large\bfseries 参考文献}
\end{center}
\vspace{2.5\baselineskip}

\VerdierBibRow{bib-001}{1}{H. Cartan 与 S. Eilenberg，“同调代数”，
  Princeton University Press (1956)。}

\VerdierBibRow{bib-002}{2}{D. Puppe，\emph{论稳定同伦论的形式结构}，
  Colloquium on Algebraic Topology, Aarhus Universitet, pp. 65-71 (1962)。}

\VerdierBibRow{bib-003}{3}{N. Yoneda，\emph{论 \textup{Ext} 与正合列}，
  J. Fac. Sci. Univ. Tokyo. Sec. I, Vol. VIII, pp. 507-576 (1960)。}

% Source-fidelity note: the printed authority reads "Harsthorne".
\VerdierBibRow{bib-004}{4}{R. Hartshorne，“留数与对偶”，Lecture Notes in
  Mathematics n° 20, Springer-Verlag (1966)。}

\VerdierBibRow{bib-005}{5}{J.-L. Verdier，\emph{概形的 étale 上同调中的一个
  对偶定理}，Lecture Notes of the Woods Hole Conference,
  American Mathematical Society (1964)。}

\VerdierBibRow{bib-006}{6}{M. Artin、A. Grothendieck、J.-L. Verdier，
  “拓扑斯理论与概形的 étale 上同调” (S.G.A. 4)，Séminaire de Géométrie
  Algébrique du Bois-Marie (1963-1964), Lecture Notes in Mathematics n° 269,
  Springer-Verlag (1972)。}

\VerdierBibRow{bib-007}{7}{J.-P. Serre，“Galois 上同调”，Lecture Notes in
  Mathematics n° 5, Springer-Verlag (1964)。}

\VerdierBibRow{bib-008}{8}{M. Artin 与 J.-L. Verdier，“数域的 étale 上同调
  讨论班”，Lecture Notes of the Woods Hole Conference,
  American Mathematical Society (1964)。}

\VerdierBibRow{bib-009}{9}{J.-L. Verdier，\emph{局部紧空间上同调中的对偶性}，
  Séminaire Bourbaki, exposé n° 300 (1965-1966)。}

\VerdierBibRow{bib-010}{10}{J.-L. Verdier 与 M. Zisman，\emph{Poincaré 对偶}，
  Séminaire Heidelberg-Strasbourg (1966-1967), Publ. I.R.M.A. Strasbourg
  n° 3 (1969)。}
\endgroup
